\section{Выводы}

В данной лабораторной работе было проведено исследование качества программы-словаря на основе дерева PATRICIA, разработанной ранее. В отличие от предыдущей лабораторной работы, здесь основное внимание уделялось не самой корректности структуры данных, а исследованию времени работы и использования памяти.

По результатам профилирования были обнаружены два явных недочёта: избыточные накладные расходы в операции сохранения дерева и слишком дорогая финальная проверка при загрузке данных из файла. Оба недочёта были исправлены и затем повторно проверены теми же средствами и на тех же входных данных.

Практический результат работы состоит в том, что удалось не только ускорить программу, но и выстроить более корректную методику исследования. Использование \texttt{gprof} позволило выделить наиболее дорогие участки по времени, а применение \texttt{Valgrind} дало возможность проверить программу с точки зрения корректности работы с памятью и оценить её потребление.

Главный вывод по лабораторной работе состоит в том, что исследование качества программы должно включать не только сами измерения, но и проверку корректности методики этих измерений. Именно это позволило перейти от общих наблюдений к конкретным улучшениям, подтверждённым повторными замерами.
\pagebreak
